\documentclass[11pt, a4paper]{article}

% ── Packages ──────────────────────────────────────────────────────────────────
\usepackage[utf8]{inputenc}
\usepackage[T1]{fontenc}
\usepackage{lmodern}
\usepackage[margin=1in]{geometry}
\usepackage{hyperref}
\usepackage{enumitem}
\usepackage{booktabs}
\usepackage{graphicx}
\usepackage{xcolor}
\usepackage{titlesec}
\usepackage{parskip}
\usepackage{fancyhdr}
\usepackage{tabularx}
\usepackage{microtype}
\usepackage{pifont}  % For checkmarks and x marks

% ── Colours & Styling ─────────────────────────────────────────────────────────
\definecolor{brand}{HTML}{2563EB}
\definecolor{accent}{HTML}{F59E0B}
\definecolor{darktext}{HTML}{1E293B}

\hypersetup{
    colorlinks=true,
    linkcolor=brand,
    urlcolor=brand,
    citecolor=brand,
}

\titleformat{\section}{\Large\bfseries\color{darktext}}{}{0em}{}[\vspace{-0.5em}\rule{\textwidth}{0.4pt}]
\titleformat{\subsection}{\large\bfseries\color{brand}}{}{0em}{}

\pagestyle{fancy}
\fancyhf{}
\fancyhead[L]{\small\textcolor{gray}{Price Sense AI}}
\fancyhead[R]{\small\textcolor{gray}{\thepage}}
\renewcommand{\headrulewidth}{0.4pt}

% ── Symbols ───────────────────────────────────────────────────────────────────
\newcommand{\cmark}{\ding{51}}  % Checkmark
\newcommand{\xmark}{\ding{55}}  % X mark

% ── Document ──────────────────────────────────────────────────────────────────
\begin{document}

% ══════════════════════════════════════════════════════════════════════════════
% COVER PAGE
% ══════════════════════════════════════════════════════════════════════════════
\begin{titlepage}
    \centering
    \vspace*{3cm}
    
    {\Huge\bfseries\textcolor{darktext}{Price Sense AI}}\\[1cm]
    {\LARGE\textcolor{brand}{Case Study Write-Up}}\\[2cm]
    {\Large AI-Powered Promotion Intelligence\\[0.3cm] for Mid-Market Retailers}\\[3cm]
    
    \begin{flushleft}
    \large
    \textbf{Submitted by:}\\[0.5cm]
    Pallab Saha\\[0.3cm]
    \href{mailto:pallabsongeet007@gmail.com}{pallabsongeet007@gmail.com}\\[0.3cm]
    +91 9062251052\\[1cm]
    
    \textbf{Challenge:}\\[0.3cm]
    24-hour prototype build + technical write-up\\[1cm]
    
    \textbf{Repository:}\\[0.3cm]
    \url{https://github.com/pallab-saha-git/Price_Sense_AI}
    \end{flushleft}
    
    \vfill
    {\large March 2026}
\end{titlepage}

\newpage

% ══════════════════════════════════════════════════════════════════════════════
\section{How Did You Decide the Features in Your POC?}
% ══════════════════════════════════════════════════════════════════════════════

The feature set was driven by a single user question: \textbf{``Should I run this promotion?''} Every feature exists only if it contributes to answering that question with data instead of gut feel.

\subsection{Discovery Process}
\begin{enumerate}[leftmargin=1.5em]
    \item \textbf{Persona mapping.} I identified three core users---Category Manager (decides which products to promote), Merchandising VP (approves the promo calendar), and Pricing/Finance Analyst (measures ROI post-hoc). Each persona has a distinct question: ``Is 25\% off better than 20\%?'', ``What's the portfolio-level impact?'', and ``What was the real incremental profit?'' respectively.

    \item \textbf{Decision decomposition.} Answering ``should I promote?'' requires five quantitative sub-answers:
    \begin{itemize}[leftmargin=1em]
        \item How much extra volume will I sell? $\rightarrow$ \textit{Price elasticity model}
        \item What would sales be without the promo? $\rightarrow$ \textit{Demand baseline forecast}
        \item Which related products will lose sales? $\rightarrow$ \textit{Cannibalization model}
        \item Will I actually make money? $\rightarrow$ \textit{Full P\&L calculator}
        \item How confident should I be? $\rightarrow$ \textit{Composite risk scorer}
    \end{itemize}
    These five models became the analytical engine. Nothing more was added to avoid scope creep.

    \item \textbf{Horizontal design constraint.} The POC must not be hardcoded to a single category. The product hierarchy was abstracted to Category $\rightarrow$ Subcategory $\rightarrow$ Brand $\rightarrow$ SKU (size/variant), which works equally for nuts, beverages, grocery/CPG, or any retail vertical.

    \item \textbf{What was deliberately excluded.} Real-time POS integration, multi-user collaboration, custom reporting, and vendor-deal management were all deferred. They are production features that add complexity without strengthening the core ``should I promote?'' answer.
\end{enumerate}

\subsection{Feature Prioritisation (MoSCoW)}

\begin{tabularx}{\textwidth}{| l | X |}
    \hline
    \textbf{Priority} & \textbf{Features} \\
    \hline\hline
    Must & Promotion analyser (5-model pipeline), recommendation card, scenario comparison, product catalog \\
    \hline
    Should & AI-generated narrative insights, dashboard KPIs, authentication \\
    \hline
    Could & Profit opportunity scanner (catalog-wide screening), async AI insight polling \\
    \hline
    Won't (now) & Real-time POS, multi-tenant, vendor deal management \\
    \hline
\end{tabularx}


% ══════════════════════════════════════════════════════════════════════════════
\section{How Did You Go About Creating the Product?}
% ══════════════════════════════════════════════════════════════════════════════

\subsection{Phase 1: Foundation \& Data Architecture}
\begin{itemize}[leftmargin=1.5em]
    \item Defined the project scope and problem decomposition: broke ``should I promote?'' into 5 quantitative sub-questions (elasticity, baseline demand, cannibalization, profit, risk).
    \item Designed the data model: 8 SQLAlchemy ORM tables covering products, stores, sales, promotions, calendar events, seasonality, customer segments, and competitor events.
    \item \textbf{Dataset selection decision:} Chose synthetic data as primary with dunnhumby as secondary real-data option (see Section~\ref{sec:dataset-choice}).
    \item \textbf{Model architecture decision:} Selected classical statistical models over deep learning (see Section~\ref{sec:model-choice}).
    \item Built the synthetic data generator producing realistic retail patterns:
    \begin{itemize}
        \item 15 SKUs (9 nuts + 6 beverages) $\times$ 25 stores $\times$ 104 weeks
        \item Embedded seasonality: Diwali $\times$3.0, Christmas $\times$2.5, summer beverages $\times$1.8
        \item Cross-elasticity coefficients: Pistachios 16oz $\leftrightarrow$ Almonds 16oz = 0.85 substitution
        \item 60 historical promotions with realistic discount depths (10--35\%)
        \item Calendar events: 15 major retail events (Black Friday, Diwali, Lunar New Year, etc.)
        \item Weather index: seasonal patterns (summer_multiplier, winter_multiplier)
        \item Customer segments: 3 personas (price-sensitive, loyalist, occasional) per store
        \item Competitor events: synthetic competitive promo activity per category
    \end{itemize}
    \item Set up SQLite database with WAL mode and seed pipeline.
\end{itemize}

\subsection{Phase 2: Analytical Engine}
\begin{itemize}[leftmargin=1.5em]
    \item Implemented all five statistical models:
    \begin{enumerate}
        \item \textbf{Elasticity:} Log-log OLS regression with seasonality (sin/cos week encoding), trend controls, and category-level fallback defaults calibrated to published retail benchmarks.
        \item \textbf{Demand Forecast:} Prophet $\rightarrow$ SARIMAX $\rightarrow$ Moving Average cascade with holiday effects and seasonality multipliers.
        \item \textbf{Cannibalization:} Cross-elasticity regression with dynamic fallbacks based on category/subcategory proximity.
        \item \textbf{Profit Calculator:} Full P\&L model computing baseline vs.~promo revenue/margin, incremental profit, ROI, with forward-buy adjustment support.
        \item \textbf{Risk Scorer:} 6-factor composite (elasticity confidence, data recency, cannibalization breadth, discount depth history, seasonal anomaly risk, category volatility).
    \end{enumerate}
    \item Built the service orchestration layer (\texttt{promo\_analyzer.py}) that calls all five models in sequence and returns a unified \texttt{PromoAnalysisResult}.
    \item Added the scenario engine for multi-discount comparison, optimised to compute elasticity/forecast once and reuse across discount levels.
\end{itemize}

\subsection{Phase 3: UI \& Intelligence Layer}
\begin{itemize}[leftmargin=1.5em]
    \item Built all Dash pages: Dashboard (KPI overview + trend charts), Analyze (promo input form + recommendation card), Compare (multi-scenario table), Catalog (product browser + promo history), and Profit Opportunities (catalog-wide screening).
    \item Created reusable UI components: recommendation card, elasticity chart, cannibalization heatmap, risk gauge, insight panel, scenario table.
    \item Implemented dual-mode insight generation:
    \begin{itemize}
        \item Template-based (deterministic, instant, always works)
        \item OpenRouter LLM integration with 8-model fallback chain, token-bucket rate limiter, circuit breakers for auth/billing failures, and async background polling
    \end{itemize}
    \item Added authentication: Flask sessions, CSRF protection (dual guard for SPA routing), bcrypt password hashing, admin seeding, registration controls.
\end{itemize}

\subsection{Phase 4: Polish, Testing \& Deployment}
\begin{itemize}[leftmargin=1.5em]
    \item Wrote unit tests for all five models, service orchestration, insight generation, auth, and data layer.
    \item Built Docker image with Prophet/CmdStan pre-compilation, Gunicorn config, non-root user, and \texttt{docker-compose.yml} for one-command deployment.
    \item Added dunnhumby real-data loader as a secondary data path (auto-detects zip files in \texttt{data/zip/}).
    \item End-to-end testing across all pages and edge cases (insufficient data, missing SKUs, extreme discounts, rate limiting, Prophet unavailability).
    \item Created comprehensive documentation (README, project outline, inline docstrings).
\end{itemize}

\subsection{AI as a Development Partner}
Throughout the build, AI (GitHub Copilot / Claude) was used as a pair-programming partner for:
\begin{itemize}[leftmargin=1.5em]
    \item Scaffolding boilerplate (ORM models, callback registration, Dockerfile)
    \item Validating statistical methodology (elasticity estimation, cross-elasticity interpretation)
    \item Debugging edge cases (Prophet installation on Windows, SQLite WAL mode, CSRF in single-page apps)
    \item Code review and refactoring suggestions
\end{itemize}
The product and architectural decisions were human-driven; AI accelerated execution.

\vspace{1em}

\subsection{Dataset Choice Rationale}\label{sec:dataset-choice}

The prototype uses \textbf{synthetic data as primary} with \textbf{dunnhumby as optional real data}. This decision was deliberate, not a shortcut.

\subsubsection{Why Synthetic Data for the Core Demo?}

\begin{enumerate}[leftmargin=1.5em]
    \item \textbf{Controllability.} Embedded cross-elasticity coefficients (e.g., Pistachios 16oz $\leftrightarrow$ Almonds 16oz = 0.85) and seasonality multipliers (Diwali $\times$3.0, summer beverages $\times$1.8) ensure models produce plausible, story-aligned outputs. Random data produces flat, implausible results.
    
    \item \textbf{Zero-barrier demo.} No download, no registration, no API keys. \texttt{python app.py} works immediately. Real datasets require manual download steps that break evaluator flow.
    
    \item \textbf{Category alignment.} Synthetic data is designed around the reference customer (specialty nuts retailer). dunnhumby has grocery departments but not explicit ``Salted Pistachios 16oz''---it requires mapping commodity codes.
    
    \item \textbf{Speed.} 15 SKUs $\times$ 25 stores $\times$ 104 weeks = 39K rows loads in $<$1 second. dunnhumby Complete Journey (2.6M transactions) takes 5--8 seconds to process.
\end{enumerate}

\subsubsection{Why dunnhumby Was Selected as Real Data Option?}

When real data \textit{is} needed (e.g., to validate model behavior on actual retail patterns), dunnhumby ``The Complete Journey'' is the only public dataset that provides all three critical features simultaneously:

\begin{tabularx}{\textwidth}{| l | c | c | c | X |}
    \hline
    \textbf{Dataset} & \textbf{Promo Flags} & \textbf{Basket Data} & \textbf{Pricing} & \textbf{Verdict} \\
    \hline\hline
    dunnhumby & \textcolor{green}{\cmark} & \textcolor{green}{\cmark} & \textcolor{green}{\cmark} & \textbf{SELECTED} --- Only dataset with all 3 signals for elasticity + cannibalization \\
    \hline
    Walmart Sales & \textcolor{orange}{\~{}} & \textcolor{red}{\xmark} & \textcolor{orange}{\~{}} & Aggregated store-week level, no SKU cross-effects visible \\
    \hline
    Corp. Favorita & \textcolor{green}{\cmark} & \textcolor{red}{\xmark} & \textcolor{green}{\cmark} & 125M rows, overkill for 24-hour build; no basket composition \\
    \hline
    Instacart & \textcolor{red}{\xmark} & \textcolor{green}{\cmark} & \textcolor{red}{\xmark} & No pricing data = cannot estimate elasticity \\
    \hline
\end{tabularx}

\textbf{The critical difference:} dunnhumby provides \texttt{basket\_id} (multiple rows = same shopping trip), which enables cannibalization detection. When a customer buys Pistachios instead of Almonds in the same trip where Pistachios are on promo, that's observable substitution. No other public dataset has this.

Additionally:
\begin{itemize}[leftmargin=1.5em]
    \item \texttt{causal\_data.csv} contains per-product per-week flags for \texttt{display} (endcap placement), \texttt{feature} (flyer ad), and \texttt{temp\_price\_reduction}---explicit promo indicators.
    \item \texttt{hh\_demographic.csv} provides household-level age, income, family size---enables customer segmentation.
    \item 2 years of data (711 days) is sufficient for seasonal pattern detection without being unwieldy.
\end{itemize}

\subsubsection{Implementation}

The app auto-detects dunnhumby zip files in \texttt{data/zip/} at startup and loads them if present. Otherwise, it falls back to synthetic data. This design ensures:
\begin{itemize}[leftmargin=1.5em]
    \item Evaluators can run the demo instantly without downloads.
    \item Technical reviewers can validate model behavior on real retail patterns by placing the zip file and restarting.
\end{itemize}

\vspace{1em}

\subsection{Model Architecture Choice}\label{sec:model-choice}

\subsubsection{Why Prophet/SARIMAX Instead of TensorFlow/PyTorch?}

The demand forecasting model uses \textbf{Prophet $\rightarrow$ SARIMAX $\rightarrow$ Moving Average cascade}, not neural networks. This was a deliberate technical decision, not a skill limitation.

\begin{tabularx}{\textwidth}{| p{3.5cm} | p{3.5cm} | X |}
    \hline
    \textbf{Criterion} & \textbf{Prophet/ SARIMAX} & \textbf{TensorFlow/ PyTorch (LSTM)} \\
    \hline\hline
    Training data & 12--52 weeks per SKU & 1000+ weeks per SKU \\
    \hline
    Interpret- ability & Explicit trend + seasonality + holidays & Black-box embeddings \\
    \hline
    Cold-start & Category fallbacks work instantly & Requires extensive training data \\
    \hline
    Inference speed & $<$100ms per SKU & 200--500ms+ (GPU needed) \\
    \hline
    Footprint & 50 MB, pure Python & 500 MB+, CUDA/cuDNN \\
    \hline
    Degradation & Prophet $\rightarrow$ SARIMAX $\rightarrow$ MA & Fails if training didn't converge \\
    \hline
    Domain fit & Designed for business forecasting & General-purpose, needs custom features \\
    \hline
    Overfitting risk & Low (regularised) & High (memorises sparse data) \\
    \hline
\end{tabularx}

\textbf{Key reasons for Prophet:}
\begin{enumerate}[leftmargin=1.5em]
    \item \textbf{Explicit seasonality + holiday modeling.} Prophet's \texttt{holidays} parameter directly accepts a DataFrame of retail events (Black Friday, Diwali, Christmas). Neural networks require manual one-hot encoding and extensive hyperparameter tuning to learn these patterns.
    
    \item \textbf{Interpretable decomposition.} Prophet outputs separate components: trend, yearly seasonality, weekly seasonality, holiday effects. A merchant can see ``+30\% from Diwali, +8\% from display effect''. LSTM outputs are opaque vectors.
    
    \item \textbf{Robust to missing data.} Retail data has gaps (stores closed, SKUs out of stock, partial promo weeks). Prophet handles missing observations natively. Neural networks require imputation strategies that introduce bias.
    
    \item \textbf{Deployment simplicity.} Prophet is pure Python (uses Stan for MCMC internally, but that's bundled). No GPU, no TensorFlow Serving, no ONNX export. Runs on any laptop or Docker container.
    
    \item \textbf{24-hour build constraint.} Training and debugging LSTM/Transformers would consume 6--8 hours (architecture search, hyperparameter tuning, convergence debugging). Prophet works out-of-the-box with sensible defaults.
\end{enumerate}

\subsubsection{Why SARIMAX as Fallback?}

SARIMAX (Seasonal ARIMA with eXogenous regressors) from \texttt{statsmodels} is the second tier in the cascade:
\begin{itemize}[leftmargin=1.5em]
    \item Zero installation --- \texttt{statsmodels} is pure Python, no binary dependencies.
    \item Works when Prophet is unavailable (Windows C++ toolchain issues).
    \item Provides seasonal decomposition (though less flexible than Prophet's additive components).
    \item Standard econometrics tool---familiar to retail analysts.
\end{itemize}

\subsubsection{When Would Neural Networks Be Appropriate?}

Deep learning becomes viable at production scale when:
\begin{itemize}[leftmargin=1.5em]
    \item \textbf{Large catalog:} 10,000+ SKUs with 3+ years of history each. Transformers can learn cross-SKU patterns.
    \item \textbf{Real-time external signals:} Weather API, competitor pricing scrapes, social media sentiment. Neural networks excel at multimodal fusion.
    \item \textbf{Accuracy gains justify complexity:} Benchmarks show $<$5\% MAPE improvement over Prophet for typical retail data. Not worth 10$\times$ deployment complexity for MVP.
\end{itemize}

For this 24-hour prototype, classical time series methods are the correct engineering choice.


% ══════════════════════════════════════════════════════════════════════════════
\section{Technical Architecture \& How It Changes with Scale}
% ══════════════════════════════════════════════════════════════════════════════

\vspace{1.5em}

\subsection{Current MVP Architecture}

{\small
\begin{verbatim}
┌──────────────────────────────────────────────────────────┐
│              Browser (Dash SPA)                          │
└────────────────────────┬─────────────────────────────────┘
                         │
                         ↓
┌──────────────────────────────────────────────────────────┐
│       Dash / Flask (app.py + Gunicorn WSGI)              │
│  ┌──────────┐  ┌────────────┐  ┌─────────────────┐      │
│  │  Pages   │  │ Callbacks  │  │  Auth (Flask    │      │
│  │(layouts) │  │  (logic)   │  │ sessions+CSRF)  │      │
│  └────┬─────┘  └─────┬──────┘  └─────────────────┘      │
│       └──────────────┼──────────────────┐                │
│                      ↓                  │                │
│  ┌───────────────────────────────────────┴─────────────┐ │
│  │     Service Layer (promo_analyzer.py)               │ │
│  │  Orchestrates: Elasticity → Forecast →              │ │
│  │  Cannibalization → P&L → Risk → Insights            │ │
│  └──────────────────────┬──────────────────────────────┘ │
│                         ↓                                │
│  ┌──────────────────────────────────────────────────┐   │
│  │   Model Layer (5 statistical models)             │   │
│  │  • Elasticity (log-log OLS + seasonality)        │   │
│  │  • Demand Forecast (Prophet/SARIMAX/MA cascade)  │   │
│  │  • Cannibalization (cross-elasticity regression) │   │
│  │  • Profit Calculator (full P&L with ROI)         │   │
│  │  • Risk Scorer (6-factor composite)              │   │
│  └──────────────────────┬───────────────────────────┘   │
│                         ↓                                │
│  ┌──────────────────────────────────────────────────┐   │
│  │   Data Layer (SQLAlchemy ORM + SQLite)           │   │
│  │  • 8 tables: products, stores, sales, promotions,│   │
│  │    calendar_events, seasonality, segments, etc.  │   │
│  │  • In-memory pandas cache (sub-2s response)      │   │
│  └──────────────────────────────────────────────────┘   │
└──────────────────────────────────────────────────────────┘
                         ↕
          ┌─────────────────────────────┐
          │ OpenRouter API (optional)   │
          │  8-model fallback chain     │
          │  for NL insights            │
          └─────────────────────────────┘
\end{verbatim}
}

\textbf{Key MVP choices:}
\begin{itemize}[leftmargin=1.5em]
    \item \textbf{SQLite} --- zero installation, file-based, sufficient for single-user demo.
    \item \textbf{In-memory pandas cache} --- all data loaded once at startup for sub-2s response times.
    \item \textbf{Synchronous model execution} --- all 5 models run in-process, sequentially.
    \item \textbf{Dash (Plotly)} --- pure Python, no React/JS required, same code MVP $\rightarrow$ prod.
\end{itemize}

\subsection{Data Features \& Intelligence Inputs}

The models leverage extensive contextual data to produce realistic, defensible recommendations:

\begin{tabularx}{\textwidth}{| l | X |}
    \hline
    \textbf{Feature Category} & \textbf{Implementation Details} \\
    \hline\hline
    Seasonality & \textbullet\, Sin/cos week-of-year encoding \newline \textbullet\, Pre-computed \texttt{seasonality\_index} table with per-SKU multipliers \newline \textbullet\, \texttt{is\_seasonal} flag + \texttt{peak\_seasons} JSON array \\
    \hline
    Calendar Events & \textbullet\, 15 retail events (Black Friday, Diwali, Christmas, etc.) \newline \textbullet\, Intensity ratings (1--5) + category tags \newline \textbullet\, Stored in \texttt{calendar\_events} table \newline \textbullet\, Used as Prophet \texttt{holidays} parameter \\
    \hline
    Weather Impact & \textbullet\, \texttt{weather\_index} table with seasonal multipliers \newline \textbullet\, Open-Meteo API integration (free, no key) \newline \textbullet\, Fallback to hardcoded averages \newline \textbullet\, Beverages: +1.8$\times$ summer lift \\
    \hline
    Customer Segments & \textbullet\, 3 personas: price-sensitive, loyalist, occasional \newline \textbullet\, Fields: \texttt{avg\_basket\_size}, \texttt{promo\_response\_multiplier}, \texttt{loyalty\_score} \newline \textbullet\, Stored per-store in \texttt{customer\_segments} \\
    \hline
    Competitor Activity & \textbullet\, \texttt{competitor\_events} table (~5 events/category/year) \newline \textbullet\, Synthetic promo dates + intensity \newline \textbullet\, Explains variance in historical promo performance \\
    \hline
    Cross-Elasticity & \textbullet\, Hardcoded substitution coefficients for within-category pairs \newline \textbullet\, Dynamic fallback: same subcat = 0.22, same cat = 0.14, cross-cat = 0.0 \newline \textbullet\, Enables realistic cannibalization heatmaps \\
    \hline
    Store Geography & \textbullet\, 25 stores across 4 regions (NE, SE, MW, W) \newline \textbullet\, Fields: \texttt{channel}, \texttt{size\_tier}, \texttt{avg\_weekly\_footfall} \newline \textbullet\, Allows region-specific elasticity estimation \\
    \hline
    Product Hierarchy & \textbullet\, Structure: Category $\rightarrow$ Subcategory $\rightarrow$ Brand $\rightarrow$ SKU \newline \textbullet\, Category-level defaults: Nuts = $-2.0$, Beverages = $-2.2$, Grocery = $-1.5$ \newline \textbullet\, Used when item-level data is sparse \\
    \hline
    Promotion History & \textbullet\, 60 historical promotions in \texttt{promotions} table \newline \textbullet\, Fields: \texttt{discount\_pct}, \texttt{promo\_type}, \texttt{units\_sold}, \texttt{revenue} \newline \textbullet\, Risk scorer checks if discount is within tested range \\
    \hline
    Margin \& Cost Data & \textbullet\, Per-SKU fields: \texttt{regular\_price}, \texttt{cost\_price}, \texttt{margin\_pct} \newline \textbullet\, Enables full P\&L calculation \newline \textbullet\, Outputs: baseline/promo margin, incremental profit, ROI \\
    \hline
\end{tabularx}

\textbf{Why this matters:} A promotion recommendation that ignores seasonality (e.g., predicting the same lift for Pistachios in July vs.~November) is implausible. Real retailers know Diwali drives 3$\times$ baseline for nuts. Embedding these patterns into synthetic data ensures the model outputs align with domain expertise.

\subsection{Production Architecture at Scale}

\begin{tabularx}{\textwidth}{| l | l | X |}
    \hline
    \textbf{Component} & \textbf{MVP} & \textbf{Production} \\
    \hline\hline
    Database        & SQLite (file)           & PostgreSQL (managed, e.g.\ AWS RDS). One config line change. \\
    \hline
    Cache           & In-process pandas       & Redis --- shared across workers, survives restarts. \\
    \hline
    Task queue      & Synchronous             & Celery + Redis --- model inference is async, non-blocking. \\
    \hline
    ML models       & Statsmodels + Prophet   & Same, plus periodic offline retraining via Airflow/Prefect. \\
    \hline
    Auth            & Flask sessions (cookie) & OAuth 2.0 / SSO (Okta, Azure AD) for enterprise. \\
    \hline
    Deployment      & Docker (single node)    & Kubernetes (EKS/GKE) with horizontal pod autoscaling. \\
    \hline
    AI insights     & OpenRouter (free tier)  & OpenAI API or self-hosted LLM behind a gateway. \\
    \hline
    Monitoring      & Loguru (stdout)         & Datadog / Prometheus + Grafana. ML model drift tracking. \\
    \hline
    CI/CD           & Manual                  & GitHub Actions $\rightarrow$ Docker build $\rightarrow$ deploy to K8s. \\
    \hline
    Multi-tenancy   & Single user             & Row-level security in Postgres, tenant isolation. \\
    \hline
\end{tabularx}

The critical insight: the \textbf{application code barely changes}. The models, service layer, and UI are the same. What changes is infrastructure wiring---database URL, cache backend, task runner, and auth provider.


% ══════════════════════════════════════════════════════════════════════════════
\section{What Challenges Did You Face?}
% ══════════════════════════════════════════════════════════════════════════════

\begin{enumerate}[leftmargin=1.5em]
    \item \textbf{Prophet installation on Windows.} Prophet requires CmdStan, which expects a C++ toolchain. The standard \texttt{pip install prophet} fails on most Windows setups. I resolved this with a Docker build that pre-installs CmdStan into Prophet's internal path, plus a model cascade (Prophet $\rightarrow$ SARIMAX $\rightarrow$ Moving Average) so the app degrades gracefully when Prophet is unavailable locally.

    \item \textbf{Making simulated data feel real.} Purely random synthetic data produces implausible elasticity estimates and flat cannibalization signals. I embedded cross-elasticity coefficients directly into the synthetic generator (e.g.\ Pistachios 16oz $\leftrightarrow$ Almonds 16oz: 0.85 substitution multiplier) and layered realistic seasonality (Diwali $\times$3.0 for nuts, summer $\times$1.8 for beverages). This ensures the models produce plausible outputs that tell a coherent story during a demo.

    \item \textbf{CSRF in a single-page app.} Dash uses HTML5 client-side routing, so clicking nav links never hits Flask's \texttt{before\_request} middleware. I had to implement a dual auth guard: server-side (Flask \texttt{before\_request} for direct URL access) plus a client-side Dash callback that monitors URL changes and forces a redirect to \texttt{/login} when the session is unauthenticated.

    \item \textbf{AI insight rate limiting.} Free-tier LLM APIs (OpenRouter) have aggressive rate limits ($\sim$8 requests/60s). Multiple rapid analyses would fail. I implemented a token-bucket rate limiter, an automatic model fallback chain (8 free models tried sequentially), circuit breakers for permanent auth/billing failures, and a temporary cooldown circuit for upstream 429 responses.

    \item \textbf{Elasticity estimation with sparse data.} Many SKUs in both synthetic and dunnhumby datasets have limited price variation, causing OLS to produce unreliable coefficients. I added category-level default elasticities (calibrated to published US grocery benchmarks from Hoch et al.\ 1995), data quality flags (\texttt{good}/\texttt{limited}/\texttt{insufficient}), and transparent communication in the UI when confidence is low.

    \item \textbf{Balancing fidelity vs.\ demo speed.} Running 5 models $\times$ 6 discount levels for scenario comparison would take 10+ seconds. I optimised the scenario engine to compute elasticity and demand forecast once (they're discount-independent) and only re-run P\&L, cannibalization, and risk per scenario, cutting response time to $\sim$2--3 seconds.
\end{enumerate}


% ══════════════════════════════════════════════════════════════════════════════
\section{What Else Would You Add Given More Time?}
% ══════════════════════════════════════════════════════════════════════════════

\begin{enumerate}[leftmargin=1.5em]
    \item \textbf{Promotion calendar view.} A Gantt-style calendar showing all planned promotions with overlap detection and portfolio-level margin impact. The VP Merchandising persona needs this to approve a full quarter's promo slate, not just one-off analyses.

    \item \textbf{Vendor deal integration.} Most retail promos are vendor-funded (the brand pays for the discount). Adding a ``vendor contribution'' field to the P\&L would flip many ``NOT RECOMMENDED'' promos to ``RECOMMENDED'' and make the tool immediately more realistic.

    \item \textbf{A/B test tracking.} Connect actual post-promo results back to the prediction. Show accuracy over time (``Our model predicted +38\% lift; actual was +41\%''). This builds trust and creates a feedback loop that improves the models.

    \item \textbf{Alerting \& notifications.} Push alerts when a competitor drops price on a category you're planning to promote, or when a planned promo's risk score shifts due to new data.

    \item \textbf{Multi-tenant SaaS architecture.} Row-level security, tenant-specific model parameters, and a self-service onboarding flow for new retailers.

    \item \textbf{Explainability layer.} SHAP-style feature attribution showing \textit{why} the model predicts a +42\% lift (e.g.\ ``30\% from price sensitivity, 8\% from seasonal demand, 4\% from display effect'').

    \item \textbf{Store-level recommendations.} Instead of chain-wide, rank which stores should run a particular promotion based on local elasticity, demographics, and competitive density.

    \item \textbf{Export \& reporting.} PDF/Excel export of analysis results for offline sharing in merchant review meetings.
\end{enumerate}


% ══════════════════════════════════════════════════════════════════════════════
\section{Anything Else Worth Sharing}
% ══════════════════════════════════════════════════════════════════════════════

\subsection{Design Philosophy}
The prototype follows a ``simulate the intelligence, not the interface'' principle. Rather than building a flashy UI over random numbers, effort was concentrated on making the \textit{analytical outputs} realistic and internally consistent:

\begin{itemize}[leftmargin=1.5em]
    \item Elasticity estimates use real statistical methods (log-log OLS with seasonality controls) grounded in published retail research.
    \item Cannibalization uses cross-elasticity regression with dynamic fallbacks based on category/subcategory proximity---not arbitrary percentages.
    \item Risk scoring is a weighted composite of six explicit factors, each with a transparent description shown in the UI.
    \item Recommendations have clear tiers (RECOMMENDED / MARGINAL / NOT\_RECOMMENDED / INSUFFICIENT\_DATA) with economic logic behind each threshold.
    \item When data is insufficient, the system says so honestly rather than showing overconfident numbers.
\end{itemize}

This means an evaluator (or a merchant in a demo) can probe any number and get a defensible explanation, which is what separates a POC from a mockup.

\subsection{Graceful Degradation}
Every external dependency has a fallback:

\begin{tabularx}{\textwidth}{| l | l | X |}
    \hline
    \textbf{Dependency} & \textbf{Primary} & \textbf{Fallback} \\
    \hline\hline
    Demand forecasting & Prophet & SARIMAX $\rightarrow$ Moving Average \\
    \hline
    AI insights        & OpenRouter LLM (8-model chain) & Deterministic templates \\
    \hline
    Dataset            & dunnhumby real data & Synthetic generator (auto) \\
    \hline
    Database           & PostgreSQL (prod) & SQLite (MVP, zero-install) \\
    \hline
\end{tabularx}

The app runs fully offline with zero API keys and zero external dependencies beyond Python packages.

\subsection{Code Quality}
\begin{itemize}[leftmargin=1.5em]
    \item Comprehensive docstrings with parameter descriptions on every module and public function.
    \item Unit tests covering all five models, service orchestration, insight generation, auth, and data layer.
    \item Structured logging (Loguru) with context-rich messages for debugging.
    \item Environment-driven configuration with production safety checks (e.g.\ blocking default \texttt{SECRET\_KEY} in non-debug mode).
\end{itemize}


% ══════════════════════════════════════════════════════════════════════════════
\section{Prototype \& Submission}
% ══════════════════════════════════════════════════════════════════════════════

\begin{tabularx}{\textwidth}{| l | X |}
    \hline
    \textbf{Artifact} & \textbf{Link / Location} \\
    \hline\hline
    Repository    & \url{https://github.com/pallab-saha-git/Price_Sense_AI} \\
    \hline
    Deployed link & \textit{Not deployed --- runs locally only} \\
    \hline
    Local run     & See instructions below \\
    \hline
\end{tabularx}

\subsection{Running the Application Locally}

\textbf{Step 1: Clone the repository}
\begin{verbatim}
git clone https://github.com/pallab-saha-git/Price_Sense_AI
cd Price_Sense_AI
\end{verbatim}

\textbf{Step 2: Run the setup script (Windows)}
\begin{verbatim}
setup.bat
\end{verbatim}
This creates a virtual environment, installs all Python dependencies, generates synthetic data, and seeds the database.

\textbf{Step 3: Start the application}
\begin{verbatim}
python app.py
\end{verbatim}
Open \texttt{http://localhost:8050} in your browser. You will be redirected to \texttt{/login}. Log in with the admin credentials set in \texttt{.env} (or use the defaults if running for the first time).

\vspace{1em}
\subsection{Running with Docker}

\textbf{Alternative: Docker Compose (all platforms)}
\begin{verbatim}
docker compose up --build
\end{verbatim}
This builds the Docker image, installs all dependencies (including Prophet/CmdStan), generates synthetic data, seeds the database, and starts the application. Open \texttt{http://localhost:8050} when ready.


% ══════════════════════════════════════════════════════════════════════════════
\section{Major Assumptions}
% ══════════════════════════════════════════════════════════════════════════════

\begin{enumerate}[leftmargin=1.5em]
    \item \textbf{Promotional elasticity $\approx$ 2$\times$ regular elasticity.} Based on Blattberg, Briesch \& Fox (1995): temporary price reductions generate 2--3$\times$ the volume response of permanent price changes, due to urgency, display effects, and reference-price anchoring.

    \item \textbf{Category-level default elasticities} are used when item-level OLS has insufficient price variation. Values are calibrated to published US grocery benchmarks (Hoch et al.\ 1995): Nuts $= -2.0$, Beverages $= -2.2$, Grocery $= -1.5$, etc.

    \item \textbf{Cross-category cannibalization is negligible.} Promoting Pistachios does not materially affect Beverage sales. Cannibalization is only modelled within-category (same category, same/adjacent subcategory).

    \item \textbf{Volume lift is capped at 80\%.} This prevents implausible forecasts from extreme elasticity $\times$ deep discount combinations. Real-world promotions rarely exceed 80\% incremental lift.

    \item \textbf{Cost prices are synthetic.} Neither the dunnhumby dataset nor the synthetic data contains real cost-of-goods. A category-specific margin percentage (typically 40--55\%) is applied to generate plausible cost prices.

    \item \textbf{Single-week promotional granularity.} Sales data is aggregated to weekly level. Daily or hourly effects (e.g.\ weekend vs.\ weekday response) are not modelled.

    \item \textbf{Forward-buy (pantry loading) is conservatively set to 0\% by default.} The model supports a forward-buy factor that reduces net incremental volume, but it is disabled unless explicitly enabled, to avoid underestimating lift in the demo.

    \item \textbf{The prototype is designed for a solo evaluator} (single-user, single-node). Multi-user concurrency, tenant isolation, and horizontal scaling are documented as production-phase changes, not POC requirements.

    \item \textbf{AI-generated insights are supplementary, not authoritative.} The LLM-generated narrative adds readability but all quantitative recommendations are driven by deterministic models. If the LLM is unavailable, the product is fully functional with template-based insights.
\end{enumerate}


\vspace{2em}
\begin{center}
    \rule{0.5\textwidth}{0.4pt}\\[6pt]
    {\small\textcolor{gray}{Price Sense AI --- Built with Dash, Python, and AI partnership.}}\\
    {\small\textcolor{gray}{Powered by Trinamix.}}
\end{center}

\end{document}
